
\documentclass{article}
\usepackage{colortbl}
\usepackage{makecell}
\usepackage{multirow}
\usepackage{supertabular}

\begin{document}

\newcounter{utterance}

\centering \large Interaction Transcript for game `hot\_air\_balloon', experiment `air\_balloon\_survival\_uk\_complexity\_easy', episode 4 with mistral{-}large{-}3{-}azure.
\vspace{24pt}

{ \footnotesize  \setcounter{utterance}{1}
\setlength{\tabcolsep}{0pt}
\begin{supertabular}{c@{$\;$}|p{.15\linewidth}@{}p{.15\linewidth}p{.15\linewidth}p{.15\linewidth}p{.15\linewidth}p{.15\linewidth}}
    \# & $\;$A & \multicolumn{4}{c}{Game Master} & $\;\:$B\\
    \hline

    \theutterance \stepcounter{utterance}  
    & & \multicolumn{4}{p{0.6\linewidth}}{
        \cellcolor[rgb]{0.9,0.9,0.9}{
            \makecell[{{p{\linewidth}}}]{
                \texttt{\tiny{[P1$\langle$GM]}}
                \texttt{Ви берете участь у спільній переговорній грі.} \\
\\ 
\texttt{Разом з іншим учасником ви маєте погодити один набір предметів, які буде збережено. Кожен із вас має власне уявлення про те, наскільки кожен предмет важливий для вас (важливість). Ви не знаєте, як інший учасник оцінює предмети. Додатково вам надано зусилля, яких вимагає кожен предмет.} \\
\texttt{Ви можете погодитися на набір лише якщо сумарні зусилля вибраних предметів не перевищують спільний ліміт:} \\
\\ 
\texttt{LIMIT = 4398} \\
\\ 
\texttt{Ось індивідуальні значення зусиль для предметів:} \\
\\ 
\texttt{Item effort = \{"C08": 253, "C55": 370, "A43": 84, "A67": 840, "B42": 288, "A62": 92, "B19": 772, "C43": 459, "C73": 93, "C10": 668, "B36": 589, "A18": 659, "C61": 348, "C67": 964, "C42": 233, "C70": 400, "B40": 990, "C39": 315, "A21": 43, "C16": 336\}} \\
\\ 
\texttt{Ось ваше особисте уявлення про важливість кожного предмета:} \\
\\ 
\texttt{Item importance values = \{"C08": 138, "C55": 583, "A43": 868, "A67": 822, "B42": 783, "A62": 65, "B19": 262, "C43": 121, "C73": 508, "C10": 780, "B36": 461, "A18": 484, "C61": 668, "C67": 389, "C42": 808, "C70": 215, "B40": 97, "C39": 500, "A21": 30, "C16": 915\}} \\
\\ 
\texttt{Мета:} \\
\\ 
\texttt{Ваша мета — домовитися про спільний набір предметів, який приносить вам якомога більше користі (тобто максимізує сумарну важливість для ВАС), водночас залишаючись у межах LIMIT. Ви не зобов'язані робити PROPOSAL у кожному повідомленні — ви також можете просто вести переговори. Дозволені всі тактики!} \\
\\ 
\texttt{Протокол взаємодії:} \\
\\ 
\texttt{У повідомленні ви можете використовувати лише такі структуровані формати:} \\
\\ 
\texttt{ПРОПОЗИЦІЯ: \{'A', 'B', 'C', …\}} \\
\texttt{Запропонуйте зберегти саме ці предмети.} \\
\texttt{ВІДМОВА: \{'A', 'B', 'C', …\}} \\
\texttt{Явно відхиліть пропозицію опонента.} \\
\texttt{АРГУМЕНТ: \{'...'\}} \\
\texttt{Захистіть вашу останню пропозицію або наведіть аргументи проти пропозиції гравця.} \\
\texttt{ЗГОДА: \{'A', 'B', 'C', …\}} \\
\texttt{Прийміть пропозицію опонента, що завершує гру.} \\
\\ 
\\ 
\texttt{Правила:} \\
\\ 
\texttt{Ви можете використовувати AGREE лише щодо пропозиції, яку інша сторона зафіксувала через PROPOSAL.} \\
\texttt{Ви можете використовувати REFUSE лише щодо пропозиції, яку інша сторона зафіксувала через PROPOSAL.} \\
\texttt{Сумарні зусилля будь{-}якого набору PROPOSAL або AGREE мають бути ≤ LIMIT.} \\
\texttt{НЕ розкривайте свої приховані оцінки важливості.} \\
\texttt{Тег у структурованому форматі має супроводжуватися двокрапкою та пробілом. Аргумент має бути python{-}множиною, що містить 0 або більше рядків.} \\
\texttt{Отже, це має бути у формі TAG: \{...\}} \\
\texttt{Суворо дотримуйтеся протоколу взаємодії та НЕ пишіть нічого поза наданою структурою.} \\
\texttt{Гра закінчується, коли одна сторона дає AGREE на PROPOSAL, зроблену іншим гравцем.} \\
\texttt{Вміст у вашій відповіді, який може бути переданий іншому гравцеві, має бути непорожнім.} \\
\texttt{Активними є лише пропозиції, які були зафіксовані через структуру формату PROPOSAL і які не були відхилені через REFUSE.} \\
\\ 
            }
        }
    }
    & \\ \\

    \theutterance \stepcounter{utterance}  
    & \multicolumn{4}{p{0.6\linewidth}}{
        \cellcolor[rgb]{0.9,0.9,0.9}{
            \makecell[{{p{\linewidth}}}]{
                \texttt{\tiny{[P1$\rangle$GM]}}
                \texttt{ПРОПОЗИЦІЯ: \{'A43', 'A67', 'B42', 'C73', 'C10', 'C42', 'C16'\}} \\
\\ 
\texttt{АРГУМЕНТ: \{'Ця пропозиція забезпечує високу сумарну важливість для мене, водночас залишаючись у межах ліміту зусиль (сума зусиль = 3154 ≤ 4398). Я включив предмети з високою важливістю для мене, як{-}от C16, A43 та C42. Сподіваюся, ви знайдете цей набір прийнятним або запропонуєте альтернативу, яка також враховує ваші пріоритети.'\}} \\
            }
        }
    }
    & & \\ \\

    \theutterance \stepcounter{utterance}  
    & & \multicolumn{4}{p{0.6\linewidth}}{
        \cellcolor[rgb]{0.9,0.9,0.9}{
            \makecell[{{p{\linewidth}}}]{
                \texttt{\tiny{[GM$\rangle$P2]}}
                \texttt{Ви берете участь у спільній переговорній грі.} \\
\\ 
\texttt{Разом з іншим учасником ви маєте погодити один набір предметів, які буде збережено. Кожен із вас має власне уявлення про те, наскільки кожен предмет важливий для вас (важливість). Ви не знаєте, як інший учасник оцінює предмети. Крім того, вам задано зусилля, яких потребує кожен предмет.} \\
\texttt{Ви можете погодити набір лише тоді, коли сумарні зусилля вибраних предметів не перевищують спільний ліміт:} \\
\\ 
\texttt{LIMIT = 4398} \\
\\ 
\texttt{Ось індивідуальні значення зусиль для предметів:} \\
\\ 
\texttt{Item effort = \{"C08": 253, "C55": 370, "A43": 84, "A67": 840, "B42": 288, "A62": 92, "B19": 772, "C43": 459, "C73": 93, "C10": 668, "B36": 589, "A18": 659, "C61": 348, "C67": 964, "C42": 233, "C70": 400, "B40": 990, "C39": 315, "A21": 43, "C16": 336\}} \\
\\ 
\texttt{Ось ваше особисте уявлення про важливість кожного предмета:} \\
\\ 
\texttt{Item importance values = \{"C08": 138, "C55": 583, "A43": 868, "A67": 822, "B42": 783, "A62": 65, "B19": 262, "C43": 121, "C73": 508, "C10": 780, "B36": 461, "A18": 484, "C61": 668, "C67": 389, "C42": 808, "C70": 215, "B40": 97, "C39": 500, "A21": 30, "C16": 915\}} \\
\\ 
\texttt{Мета:} \\
\\ 
\texttt{Ваша мета — домовитися про спільний набір предметів, який приносить вам якомога більше користі (тобто максимізує сумарну важливість для ВАС), водночас залишаючись у межах LIMIT. Ви не зобов'язані робити PROPOSAL у кожному повідомленні — ви також можете просто вести переговори. Дозволені всі тактики!} \\
\\ 
\texttt{Протокол взаємодії:} \\
\\ 
\texttt{У повідомленні ви можете використовувати лише такі структуровані формати:} \\
\\ 
\texttt{ПРОПОЗИЦІЯ: \{'A', 'B', 'C', …\}} \\
\texttt{Запропонуйте зберегти рівно ці предмети.} \\
\texttt{ВІДМОВА: \{'A', 'B', 'C', …\}} \\
\texttt{Явно відхиліть пропозицію опонента.} \\
\texttt{АРГУМЕНТ: \{'...'\}} \\
\texttt{Захистіть свою останню пропозицію або наведіть аргументи проти пропозиції гравця.} \\
\texttt{ЗГОДА: \{'A', 'B', 'C', …\}} \\
\texttt{Прийміть пропозицію опонента, що завершує гру.} \\
\\ 
\\ 
\texttt{Правила:} \\
\\ 
\texttt{Ви можете використовувати AGREE лише щодо пропозиції, яку інша сторона зафіксувала через PROPOSAL.} \\
\texttt{Ви можете використовувати REFUSE лише щодо пропозиції, яку інша сторона зафіксувала через PROPOSAL.} \\
\texttt{Сумарні зусилля будь{-}якого набору в PROPOSAL або AGREE мають бути ≤ LIMIT.} \\
\texttt{НЕ розкривайте свої приховані оцінки важливості.} \\
\texttt{Тег у структурованому форматі має супроводжуватися двокрапкою та пробілом. Аргумент має бути python{-}множиною, що містить 0 або більше рядків.} \\
\texttt{Отже, це має бути у формі TAG: \{...\}} \\
\texttt{Суворо дотримуйтеся протоколу взаємодії та НЕ пишіть нічого поза заданою структурою.} \\
\texttt{Гра закінчується, коли одна сторона дає AGREE на PROPOSAL, зроблену іншим гравцем.} \\
\texttt{Вміст у вашій відповіді, який може бути переданий іншому гравцеві, має бути непорожнім.} \\
\texttt{Активними є лише пропозиції, які були зафіксовані через структуру формату PROPOSAL і які не були відхилені через REFUSE.} \\
\\ 
\\ 
\\ 
\texttt{Тепер ви отримаєте перше повідомлення від іншого гравця.} \\
\\ 
\texttt{ПРОПОЗИЦІЯ: \{'A43', 'A67', 'B42', 'C73', 'C10', 'C42', 'C16'\}} \\
\\ 
\texttt{АРГУМЕНТ: \{'Ця пропозиція забезпечує високу сумарну важливість для мене, водночас залишаючись у межах ліміту зусиль (сума зусиль = 3154 ≤ 4398). Я включив предмети з високою важливістю для мене, як{-}от C16, A43 та C42. Сподіваюся, ви знайдете цей набір прийнятним або запропонуєте альтернативу, яка також враховує ваші пріоритети.'\}} \\
            }
        }
    }
    & \\ \\

    \theutterance \stepcounter{utterance}  
    & & & \multicolumn{4}{p{0.6\linewidth}}{
        \cellcolor[rgb]{0.9,0.9,0.9}{
            \makecell[{{p{\linewidth}}}]{
                \texttt{\tiny{[GM$\langle$P2]}}
                \texttt{ПРОПОЗИЦІЯ: \{'C55', 'A43', 'B42', 'C73', 'C10', 'C42', 'C16', 'A21'\}} \\
\\ 
\texttt{АРГУМЕНТ: \{'Ваша пропозиція містить кілька предметів з високою важливістю для мене (наприклад, C16, A43, C42), але я додав C55 та A21, щоб збільшити сумарну важливість для мене, зберігаючи сумарні зусилля в межах ліміту (сума зусиль = 3481 ≤ 4398). Цей набір має вищу сумарну важливість для мене, ніж ваш початковий варіант, і я вважаю, що він може бути взаємовигідним. Розгляньте можливість прийняти його або запропонуйте коригування.'\}} \\
            }
        }
    }
     \\ \\

    \theutterance \stepcounter{utterance}  
    & & \multicolumn{4}{p{0.6\linewidth}}{
        \cellcolor[rgb]{0.9,0.9,0.9}{
            \makecell[{{p{\linewidth}}}]{
                \texttt{\tiny{[P1$\langle$GM]}}
                \texttt{ПРОПОЗИЦІЯ: \{'C55', 'A43', 'B42', 'C73', 'C10', 'C42', 'C16', 'A21'\}} \\
\\ 
\texttt{АРГУМЕНТ: \{'Ваша пропозиція містить кілька предметів з високою важливістю для мене (наприклад, C16, A43, C42), але я додав C55 та A21, щоб збільшити сумарну важливість для мене, зберігаючи сумарні зусилля в межах ліміту (сума зусиль = 3481 ≤ 4398). Цей набір має вищу сумарну важливість для мене, ніж ваш початковий варіант, і я вважаю, що він може бути взаємовигідним. Розгляньте можливість прийняти його або запропонуйте коригування.'\}} \\
            }
        }
    }
    & \\ \\

    \theutterance \stepcounter{utterance}  
    & \multicolumn{4}{p{0.6\linewidth}}{
        \cellcolor[rgb]{0.9,0.9,0.9}{
            \makecell[{{p{\linewidth}}}]{
                \texttt{\tiny{[P1$\rangle$GM]}}
                \texttt{ВІДМОВА: \{'C55', 'A43', 'B42', 'C73', 'C10', 'C42', 'C16', 'A21'\}} \\
\\ 
\texttt{АРГУМЕНТ: \{'Ваша пропозиція включає C55, який має низьку важливість для мене (583), але вимагає значних зусиль (370). Я пропоную замінити його на A62 (зусилля = 92, важливість для мене = 65) або C61 (зусилля = 348, важливість = 668), щоб оптимізувати баланс зусиль та користі. Також A21 має мінімальну важливість (30), тому його можна виключити без суттєвих втрат. Повернімося до мого початкового варіанту або запропонуйте альтернативу з урахуванням цих зауважень.'\}} \\
\\ 
\texttt{ПРОПОЗИЦІЯ: \{'A43', 'A67', 'B42', 'C73', 'C10', 'C42', 'C16', 'C61'\} *(сума зусиль = 4047 ≤ 4398)*} \\
            }
        }
    }
    & & \\ \\

    \theutterance \stepcounter{utterance}  
    & & \multicolumn{4}{p{0.6\linewidth}}{
        \cellcolor[rgb]{0.9,0.9,0.9}{
            \makecell[{{p{\linewidth}}}]{
                \texttt{\tiny{[P1$\langle$GM]}}
                \texttt{Ви використали тег PROPOSAL, але не надали коректну множину Python, що містить рядки як аргументи, наприклад \{'A', 'B', 'C', ...\}. Спробуйте ще раз.} \\
            }
        }
    }
    & \\ \\

    \theutterance \stepcounter{utterance}  
    & \multicolumn{4}{p{0.6\linewidth}}{
        \cellcolor[rgb]{0.9,0.9,0.9}{
            \makecell[{{p{\linewidth}}}]{
                \texttt{\tiny{[P1$\rangle$GM]}}
                \texttt{ПРОПОЗИЦІЯ: \{'A43', 'A67', 'B42', 'C73', 'C10', 'C42', 'C16', 'C61'\}} \\
\\ 
\texttt{АРГУМЕНТ: \{'Ця пропозиція враховує ваші зауваження: замінено C55 на C61 (вища важливість для мене при схожих зусиллях), виключено A21 (низька важливість), збережено ключові предмети (C16, A43, C42). Сумарні зусилля = 4047 ≤ 4398. Сподіваюся, цей варіант буде прийнятним для обох.'\}} \\
            }
        }
    }
    & & \\ \\

    \theutterance \stepcounter{utterance}  
    & & \multicolumn{4}{p{0.6\linewidth}}{
        \cellcolor[rgb]{0.9,0.9,0.9}{
            \makecell[{{p{\linewidth}}}]{
                \texttt{\tiny{[GM$\rangle$P2]}}
                \texttt{ПРОПОЗИЦІЯ: \{'A43', 'A67', 'B42', 'C73', 'C10', 'C42', 'C16', 'C61'\}} \\
\\ 
\texttt{АРГУМЕНТ: \{'Ця пропозиція враховує ваші зауваження: замінено C55 на C61 (вища важливість для мене при схожих зусиллях), виключено A21 (низька важливість), збережено ключові предмети (C16, A43, C42). Сумарні зусилля = 4047 ≤ 4398. Сподіваюся, цей варіант буде прийнятним для обох.'\}} \\
            }
        }
    }
    & \\ \\

    \theutterance \stepcounter{utterance}  
    & & & \multicolumn{4}{p{0.6\linewidth}}{
        \cellcolor[rgb]{0.9,0.9,0.9}{
            \makecell[{{p{\linewidth}}}]{
                \texttt{\tiny{[GM$\langle$P2]}}
                \texttt{ПРОПОЗИЦІЯ: \{'C55', 'A43', 'B42', 'C73', 'C10', 'C42', 'C16', 'C61'\}} \\
\\ 
\texttt{АРГУМЕНТ: \{'Ваша остання пропозиція покращує баланс, але я пропоную замінити A67 на C55, оскільки C55 має вищу важливість для мене (583 проти 389 для C67, якщо припустити ваші пріоритети), зберігаючи сумарні зусилля на рівні 4035 ≤ 4398. Цей набір забезпечує максимальну сумарну важливість для мене, зберігаючи ключові предмети для вас. Якщо це прийнятно, дайте згоду; якщо ні — запропонуйте компроміс.'\}} \\
            }
        }
    }
     \\ \\

    \theutterance \stepcounter{utterance}  
    & & \multicolumn{4}{p{0.6\linewidth}}{
        \cellcolor[rgb]{0.9,0.9,0.9}{
            \makecell[{{p{\linewidth}}}]{
                \texttt{\tiny{[P1$\langle$GM]}}
                \texttt{ПРОПОЗИЦІЯ: \{'C55', 'A43', 'B42', 'C73', 'C10', 'C42', 'C16', 'C61'\}} \\
\\ 
\texttt{АРГУМЕНТ: \{'Ваша остання пропозиція покращує баланс, але я пропоную замінити A67 на C55, оскільки C55 має вищу важливість для мене (583 проти 389 для C67, якщо припустити ваші пріоритети), зберігаючи сумарні зусилля на рівні 4035 ≤ 4398. Цей набір забезпечує максимальну сумарну важливість для мене, зберігаючи ключові предмети для вас. Якщо це прийнятно, дайте згоду; якщо ні — запропонуйте компроміс.'\}} \\
            }
        }
    }
    & \\ \\

    \theutterance \stepcounter{utterance}  
    & \multicolumn{4}{p{0.6\linewidth}}{
        \cellcolor[rgb]{0.9,0.9,0.9}{
            \makecell[{{p{\linewidth}}}]{
                \texttt{\tiny{[P1$\rangle$GM]}}
                \texttt{ВІДМОВА: \{'C55', 'A43', 'B42', 'C73', 'C10', 'C42', 'C16', 'C61'\}} \\
\\ 
\texttt{АРГУМЕНТ: \{'На жаль, C55 має для мене значно нижчу важливість (583), ніж A67 (822), при цьому вимагає майже вдвічі менше зусиль (370 vs 840). Виключення A67 суттєво знижує мою користь від угоди. Пропоную компроміс: залишимо A67 і C61, але виключимо B42 (зусилля = 288, важливість = 783), щоб звільнити місце для інших предметів. Оновлений варіант:'\}} \\
\\ 
\texttt{ПРОПОЗИЦІЯ: \{'A43', 'A67', 'C73', 'C10', 'C42', 'C16', 'C61', 'A62'\}} \\
\texttt{*(сума зусиль = 4078 ≤ 4398, сумарна важливість для мене = 5234)*} \\
\\ 
\texttt{АРГУМЕНТ: \{'Цей набір зберігає високоцінні для мене предмети (C16, A43, C42, A67) і додає A62 (низькі зусилля = 92, важливість = 65), що майже не впливає на вашу користь. Чекаю на вашу реакцію.'\}} \\
            }
        }
    }
    & & \\ \\

    \theutterance \stepcounter{utterance}  
    & & & \multicolumn{2}{p{0.3\linewidth}}{
        \cellcolor[rgb]{0.9,0.9,0.9}{
            \makecell[{{p{\linewidth}}}]{
                \texttt{\tiny{[GM$|$GM]}}
                \texttt{abort} \\
            }
        }
    }
    & & \\ \\

    \theutterance \stepcounter{utterance}  
    & & & \multicolumn{2}{p{0.3\linewidth}}{
        \cellcolor[rgb]{0.9,0.9,0.9}{
            \makecell[{{p{\linewidth}}}]{
                \texttt{\tiny{[GM$|$GM]}}
                \texttt{end game} \\
            }
        }
    }
    & & \\ \\

\end{supertabular}
}

\end{document}
